% Stage A, Stage B and the Stage C feasibility gate: results report.
%
% Formatting. This file targets the ICLR conference template. The official style file is
% distributed by the conference and is not on CTAN, so it cannot be fetched automatically.
% Drop iclr2026_conference.sty (and iclr2026_conference.bst) beside this file and the
% document switches to it with no edit; without it the fallback preamble below approximates
% the same layout using standard packages, so the report compiles today either way.
%
%   latexmk -pdf docs/results_report.tex

\documentclass{article}

\IfFileExists{iclr2026_conference.sty}{%
  \usepackage{iclr2026_conference}%
  \usepackage{times}%
  \newcommand{\fallbacknote}{}%
}{%
  \usepackage[letterpaper,margin=1in]{geometry}%
  \usepackage{times}%
  \usepackage{titlesec}%
  \titleformat{\section}{\normalfont\large\bfseries}{\thesection}{1em}{\uppercase}%
  \titleformat{\subsection}{\normalfont\normalsize\bfseries}{\thesubsection}{1em}{}%
  \newcommand{\fallbacknote}{%
    \begin{center}\small\itshape
    Typeset with the fallback preamble: \texttt{iclr2026\_conference.sty} was not found.
    \end{center}}%
}

\usepackage{amsmath,amssymb}
\usepackage{booktabs}
\usepackage{graphicx}
\usepackage{caption}
\usepackage[hidelinks]{hyperref}
\usepackage{url}

\graphicspath{{figures/}{../handoff/results/}}

\title{Task Coordinates for Few-Shot Diffusion Adaptation:\\
Source Dependence, Data Scarcity and Map Structure}

\author{Anonymous authors\\Paper under double-blind review}

\date{17 September 2026}

\begin{document}

\maketitle
\fallbacknote

\begin{abstract}
We report the results of the current experimental stage for Meta-Diffusion, a method that
adapts a frozen denoising diffusion probabilistic model to a new data domain by fitting a
low-dimensional task coordinate rather than by updating network weights. The stage was
governed by a criterion fixed in advance: the main claim is supported only if knowledge of
the correct source task, combined with a learned source-to-target relation, outperforms
\emph{both} transfer that knows only the relation \emph{and} few-shot modelling that uses
only target-domain data. On the distributional metric named as headline by the plan, that
criterion is met in the low-target-data regime and only there. With a single target image
the method improves on relation-only transfer by $+0.025 \pm 0.008$ and on target-only
modelling by $+0.018 \pm 0.009$; the advantage over target-only modelling has vanished by
twenty target images. Three supporting results accompany it: the mechanism behaves as
designed on an analytically tractable task family, the marginal value of source evidence is
exhausted below the smallest budget measured, and the transport map must exceed the
identity but need not be non-linear. Two earlier readings are withdrawn, both traceable to
a measurement instrument that was blind to the quantity under test.
\end{abstract}

\section{Objective and the criterion for the claim}

The method under test factors a generative model into two parts of very unequal dimension:
a high-dimensional score network whose weights are shared across all tasks and frozen at
meta-test time, and a $k$-dimensional task coordinate that carries everything specific to a
task. Adaptation to a new domain is then the problem of choosing one short vector rather
than of retraining a network. A transport map predicts the coordinate for a target domain
from source-domain data alone, which is what makes the method applicable when target data
are scarce.

The plan governing this stage states the criterion for the main claim without
qualification: the claim is supported only if the correct source task \emph{plus} the
learned source-to-target relation outperforms \emph{both} relation-only transfer \emph{and}
target-only few-shot modelling, in the low-target-data regime. Any other outcome is to be
reported as what it is. This report is organised around that sentence.

\section{Method, instruments and terminology}

\paragraph{The model.} The backbone is a denoising diffusion probabilistic model: a network
trained to predict the noise added to an image at a randomly chosen point along a fixed
corruption schedule, so that running the schedule backwards generates samples. Generation
in this report uses the denoising diffusion implicit model sampler with the stochasticity
parameter set to zero, which makes sampling deterministic given the initial noise.

\paragraph{The coordinate and the transport map.} A set encoder $r_\psi$ maps a set of
images to a task coordinate $z \in \mathbb{R}^{k}$. Writing $z_S$ for the coordinate
encoded from source-domain data and $c$ for a descriptor of the source-to-target relation,
the transport map produces the predicted target coordinate
$\tilde{z}_T = z_S + \Delta_\gamma(z_S, c)$. Critically, $\Delta_\gamma$ takes no target
image as input; target images enter only through an optional refinement step that minimises
the denoising loss over the coordinate alone, with the network frozen.

\paragraph{Evidence budgets.} $M_S$ denotes the number of source-domain images the encoder
sees when producing $z_S$. $K_T$ denotes the number of target-domain images available at
adaptation time. These are the two axes swept in this stage.

\paragraph{The strategies compared.} \emph{No adaptation} sets the coordinate to zero.
\emph{Reuse $z_S$} transfers the source coordinate unchanged. \emph{Transport} applies
$\Delta_\gamma$; \emph{transport + refine} then refines on the $K_T$ target images.
\emph{Target only} encodes those same $K_T$ images directly and refines from there, using no
source data at all, and is therefore the hardest baseline: beating it is what demonstrates
that source data are worth having. \emph{Oracle} refines on abundant target data and bounds
what the coordinate can achieve within this basis and backbone.

\paragraph{The source controls.} Three conditions destroy task identity while preserving the
relation. \emph{Mean $z_S$} substitutes the average source coordinate over all episodes
sharing a transformation. \emph{Shuffled $z_S$} substitutes another class's coordinate drawn
from within the same transformation. \emph{Relation-only} evaluates $\Delta_\gamma(0, c)$,
keeping the relation descriptor and destroying the coordinate entirely.

\paragraph{Two instruments, and why both are reported.} The first is a transformation
statistic: three fixed convolutions measuring mean absolute Laplacian response, per-image
standard deviation, and residual above a local average. It reads which image transformation
was applied and nothing else. The second is the sliced Wasserstein distance in pixel space,
which compares two sets of images by projecting both onto many random directions, sorting
each projection, and comparing quantiles; it is the distributional metric the plan names as
headline. Section~\ref{sec:a1} shows that the choice between them is not cosmetic.

\paragraph{Statistics.} Every comparison is a paired difference computed episode by episode,
reported with a 95\% confidence interval; an interval that excludes zero is marked with an
asterisk. An interval containing zero means the sign of the difference is not determined at
this sample size; it does not establish equality, which would require a separately declared
margin of indifference that none of these comparisons carries.

\section{Experimental configuration}

All results in Sections~\ref{sec:a1} through \ref{sec:a5} use one checkpoint of the
three-transformation CIFAR-100 domain-shift task family, trained for 50{,}000 steps, read
with exponentially moving averaged weights. The task family is Gaussian blur, additive
Gaussian noise and contrast reduction at severity 3; episodes pair a clean source domain
with a corrupted target domain of the same semantic class.

\begin{table}[t]
\centering
\caption{Training and adaptation configuration. Values are read from the checkpoint's own
stored configuration and from the adaptation budget object shared by every strategy.}
\label{tab:config}
\small
\begin{tabular}{llll}
\toprule
\multicolumn{2}{l}{\textbf{Training}} & \multicolumn{2}{l}{\textbf{Adaptation (meta-test)}}\\
\midrule
Backbone & U-Net, 128 base channels & Refinement steps $J$ & 25\\
Channel multipliers & $(1,2,2)$, 2 residual blocks & Coordinate rate $\eta_z$ & $10^{-2}$\\
Coordinate dimension $k$ & 16 & Prior weight $\beta_0$ & 1.0, divided by $K_T$\\
Set encoder pooling & mean & Noise draws per step & 32\\
Transport hidden width & 64 & & \\
Relation descriptor $d_c$ & 8 & \multicolumn{2}{l}{\textbf{Sampling}}\\
Diffusion steps & 1000, cosine schedule & Sampler & implicit, $\eta=0$\\
Optimiser & AdamW, $2\times10^{-4}$ & Reverse steps & 50\\
Warmup, gradient clip & 500 steps, 1.0 & Samples per cell & 96\\
Moving-average decay & 0.999 & Episodes per cell & 60\\
Loss weights $\lambda_T, \lambda_{tr}$ & 1.0, 1.0 & & \\
Coordinate penalty $\lambda_z$ & $10^{-4}$ & & \\
Source budget at training $M_S$ & 64 & & \\
\bottomrule
\end{tabular}
\end{table}

\paragraph{A departure from the written configuration, now corrected.} The configuration
file recorded 32 noise draws per refinement step only after this stage began. The field
previously read 16 and was never passed to the object that consumes it, so every number
produced before the correction was in fact generated with that object's own default of 32.
The field was changed to match what had been running rather than the reverse. Any earlier
document quoting 16 is wrong.

\paragraph{A change that did alter conclusions.} Refinement originally drew its noise from
an unseeded global stream, so two strategies compared within one episode refined against
different noise: an unpaired difference inside an otherwise paired comparison. A per-episode
stream, rewound for each strategy, was introduced part-way through the stage. Two readings
moved when it was, and both are identified where they occur.

\section{Stage B: the positive control}

Before interpreting a small effect on images, it is worth establishing that the mechanism
works at all where the ground truth is analytic. Stage B applies the same construction to
two-dimensional Gaussian mixtures across five task families spanning a range of
source-to-target distances, with 32 tasks each.

The mechanism behaves exactly as designed. The paired improvement from using a task's own
coordinate rather than the coordinate shared by all tasks runs from $+0.1428$ to $+0.1571$
across the five families, every interval excluding zero, and the task-specific share of the
gain is $1.0$: on analytic mixtures essentially the whole benefit of the coordinate is
specific to the task rather than a constant offset. The coordinate-loss profile is minimised
at the task's own coordinate in all five families, which is the correctness check the panel
exists to perform.

Figure~\ref{fig:profile} places this beside the corresponding measurement on images. The two
loss scales are not comparable, so the figure plots the depth of the well as a percentage of
the loss at the shared mean coordinate, which is. The well is 26.1\% deep on mixtures and
1.6\% deep on CIFAR-100, and both are minimised in the right place with the wrong direction
rising monotonically. The mechanism on images is therefore neither absent nor misdirected.
It is shallow, and an earlier reading that described it as flat was overstating.

\begin{figure}[t]
\centering
\includegraphics[width=\textwidth]{profile_contrast.png}
\caption{The coordinate-loss profile on two-dimensional Gaussian mixtures (left) and on
CIFAR-100 domain shift (right). The horizontal axis walks the coordinate from the shared
mean at $s=0$ through the task's own coordinate at $s=1$ and out to twice that displacement.
Both panels are minimised at $s=1$; the depth of the well differs by roughly sixteenfold.}
\label{fig:profile}
\end{figure}

\paragraph{One result the sweep did not produce.} The relatedness sweep was built to vary
the distance between source and target tasks, and it does not move the task-specific term:
the five readings are $+0.1539$, $+0.1428$, $+0.1536$, $+0.1571$ and $+0.1533$, from the
smallest perturbation through to unrelated task pairs. The dial turns; the quantity it was
built to move does not. We record this as an open question rather than a footnote.

\section{A1: source-dependence controls}
\label{sec:a1}

This panel decides the main claim. It compares the full method against the three ways of
destroying task identity while preserving the relation, over 180 episodes at
$K_T \in \{1, 5, 20\}$.

\begin{table}[t]
\centering
\caption{A1, paired differences on the headline metric. Positive values favour the full
method. All nine cells exclude zero. The final row reports the same comparison against mean
$z_S$ under the transformation statistic, where every interval contains zero.}
\label{tab:a1}
\small
\begin{tabular}{lccc}
\toprule
Full method against & $K_T=1$ & $K_T=5$ & $K_T=20$\\
\midrule
Mean $z_S$ & $+0.0301 \pm 0.0072^{*}$ & $+0.0346 \pm 0.0085^{*}$ & $+0.0338 \pm 0.0080^{*}$\\
Shuffled $z_S$ & $+0.0416 \pm 0.0125^{*}$ & $+0.0453 \pm 0.0149^{*}$ & $+0.0436 \pm 0.0144^{*}$\\
Relation-only & $+0.0254 \pm 0.0078^{*}$ & $+0.0285 \pm 0.0088^{*}$ & $+0.0295 \pm 0.0087^{*}$\\
Target-only & $+0.018 \pm 0.009^{*}$ & $+0.011 \pm 0.007^{*}$ & $-0.003 \pm 0.007$\\
\midrule
Mean $z_S$, transformation statistic & $+0.017 \pm 0.033$ & $-0.008 \pm 0.030$ & $+0.011 \pm 0.028$\\
\bottomrule
\end{tabular}
\end{table}

\paragraph{The criterion is met, with a boundary.} At a single target image the method beats
relation-only transfer and target-only modelling simultaneously, which is what the plan
requires. The advantage over target-only modelling survives to five target images and is
gone by twenty. The claim this panel supports is therefore explicitly bounded: source
information is worth having when target evidence is scarce, and its value decays as target
evidence accumulates. That boundary is consistent with the task family, in which a single
target image already reveals which transformation was applied.

\paragraph{The two instruments disagree, and the disagreement is itself the finding.} Under
the transformation statistic the comparison against mean $z_S$ contains zero at every $K_T$.
That instrument reads which transformation was applied; the within-relation mean preserves
the transformation exactly and destroys only task identity. The instrument is therefore
blind by construction to precisely the quantity the control removes, and a null reading from
it carries no information about that quantity. An earlier reading of this panel concluded
that the task-specific content of the coordinate was unused. That conclusion was drawn from
the blind instrument and is withdrawn.

\paragraph{The gap to the oracle is real.} On the headline metric the full method remains
behind the oracle by $+0.026$, $+0.027$ and $+0.017$ at the three budgets, every interval
excluding zero. Earlier statements that the method was indistinguishable from the oracle came
from the loss-shaped instrument and are likewise withdrawn.

\paragraph{Refinement is not reportable.} The two instruments disagree on its sign, and four
measurements of the $K_T=1$ cell span 0.023 while their half-widths are approximately 0.014.
Following the plan's own exclusion list, refinement stays out of the headline method. Its
value is real when the starting coordinate is far away and unmeasurable once transport has
landed close.

\section{A3: target-data scarcity}

Over 300 episodes at $K_T \in \{1,2,5,10,20\}$, the method beats reusing the source
coordinate at every budget, the advantage decaying monotonically from
$+1.140 \pm 0.108$ to $+0.229 \pm 0.059$ with every interval clear of zero. This is the
largest and most robust effect in the stage, an order of magnitude above every other
comparison, and it establishes that the transport map does substantive work.

\paragraph{A retraction.} An earlier run had read the twenty-image budget as the method
\emph{losing} to target-only modelling, $-0.017 \pm 0.044$, and this was written up as a
crossing point. The larger run reads $+0.008 \pm 0.036$: the sign reversed and both
intervals contain zero, so the method was never shown to lose. The correct statement is that
the two become indistinguishable, not that they cross.

\section{A4: source abundance}

\begin{figure}[t]
\centering
\includegraphics[width=\textwidth]{A4_ms_curve.png}
\caption{A4. Six source budgets from 16 to 300 images, one target image, 60 episodes each,
one checkpoint throughout. Left: absolute distance per strategy. Right: paired differences
with 95\% intervals. Nothing moves.}
\label{fig:a4}
\end{figure}

The plan anticipated improvement with increasing source evidence followed by saturation.
What the sweep shows is saturation from the first point measured. Across a nineteenfold
range the comparison against reusing the source coordinate moves 0.0002, that against
target-only modelling 0.0019, and that against the oracle 0.0011, each smaller than the
half-width of its own interval at either end. Sixteen source images already buy everything
this transport map extracts from a source set.

Two strategies in the panel never consult the source coordinate, and they return identical
values at all six budgets. That is the sweep's internal check that it varies only what it
claims to vary, and it passes.

\paragraph{Scope.} Every point of this sweep is the same checkpoint, trained once at
$M_S=64$, evaluated with a different number of source images. The sweep therefore measures
how much source evidence a trained transport map requires at deployment. It does not
measure how that map would differ had it been trained under a different source budget; that
is a distinct experiment, discussed in Section~\ref{sec:deferred}.

\section{A5: transport-map structure}
\label{sec:a5}

Four maps were trained separately under identical conditions and compared episode by
episode: the residual multilayer perceptron used by the method (6824 parameters), the
identity (0), a constant relation vector (96), and a linear map (352).

\begin{table}[t]
\centering
\caption{A5, paired against the residual multilayer perceptron. Positive values favour it.
The final row gives the noise floor: the spread across the same four runs of a strategy that
never consults the transport map, and therefore reflects ordinary run-to-run variation.}
\label{tab:a5}
\small
\begin{tabular}{lccc}
\toprule
Variant & $K_T=1$ & $K_T=5$ & $K_T=20$\\
\midrule
Identity & $+0.0242 \pm 0.0080^{*}$ & $+0.0256 \pm 0.0089^{*}$ & $+0.0246 \pm 0.0087^{*}$\\
Constant relation vector & $+0.0082 \pm 0.0051^{*}$ & $+0.0092 \pm 0.0052^{*}$ & $+0.0109 \pm 0.0055^{*}$\\
Linear map & $+0.0022 \pm 0.0049$ & $+0.0004 \pm 0.0048$ & $+0.0036 \pm 0.0050$\\
\midrule
Run-to-run floor & 0.0076 & 0.0036 & 0.0135\\
\bottomrule
\end{tabular}
\end{table}

These are four separately trained models, so a difference between them mixes the structure
of the map with ordinary variation between training runs. The panel therefore prints its own
floor: a strategy that never consults the transport map spreads by 0.0076 to 0.0135 across
the same four runs, against a transport spread of 0.0082 to 0.0109 once the identity is
excluded.

Read against that floor, the defensible conclusion is narrow. The map must be more than the
identity, by a margin comfortably above the floor at every budget. The constant relation
vector's margin sits inside the floor. The linear map is indistinguishable from the residual
multilayer perceptron at all three budgets, with nineteen times fewer parameters. On these
numbers the linear map is the structure to prefer, and the claim that the non-linear map
earns its place is not supported.

\section{C0: feasibility of the demographic experiment}

The demographic scarcity experiment requires a dermatological image dataset annotated by
skin phototype. The annotation table is fully available: 16{,}577 rows over 114 conditions.
The images are not. One of the two source archives holds 12{,}631 of those rows, 76.2\% of
the dataset, and its domain fails to resolve for public resolvers as well as locally, with
the delegation itself reported as unreachable; a sample of thirty of its links returned no
image bytes at all. The second archive is fully live, and a sample of thirty of its links
returned thirty images.

Grouping phototypes I--III against IV--VI, chosen from the observed counts rather than in
advance, the reachable archive yields 16 conditions above the plan's floor. Only 8 of those
can fund a twenty-image support set together with a held-out query set, which gives a
five/two/one split of tasks, and a single held-out test task cannot carry a paired interval.
The reachable subset is therefore scaffolding on which the domain module and evaluation
predictors can be built, not a basis for the headline scarcity result. Restoring the full
dataset requires a request form that the authors control.

\section{Deferred experiments and the reasons}
\label{sec:deferred}

\begin{table}[t]
\centering
\caption{Status of the ten deliverables fixed by the plan.}
\label{tab:status}
\small
\begin{tabular}{llp{7.2cm}}
\toprule
\# & Deliverable & Status\\
\midrule
1 & Gaussian-mixture mechanism result & Complete\\
2 & Source-dependence controls & Complete\\
3 & Additive basis against feature-wise linear modulation & \textbf{Deferred}\\
4 & Target-scarcity curve & Complete\\
5 & Source-abundance curve & Complete\\
6 & Transport-structure ablation & Complete\\
7 & One demographic scarcity experiment & \textbf{Deferred}\\
8 & Paired 95\% intervals for every headline comparison & Complete\\
9 & Downstream demographic augmentation (optional) & \textbf{Deferred}\\
10 & Few-shot face generation stress test (optional) & \textbf{Deferred}\\
\bottomrule
\end{tabular}
\end{table}

\paragraph{Deliverable 3, the comparison against feature-wise linear modulation.} Feature-wise
linear modulation is the standard alternative to an additive basis: instead of adding a
learned direction to the score, it scales and shifts intermediate feature maps as a function
of the task coordinate. The comparison is deferred because of a repository state, not a
scientific judgement. The implementation of that arm and the generation scripts that would
evaluate it live on different branches, and the integration step that the plan places before
all of Stage A has not been performed. Until those branches are merged no cross-arm
measurement can run. The arm itself is built and has been measured on the denoising loss;
only the generation-level comparison is outstanding.

\paragraph{Deliverables 7 and 9, the demographic experiments.} The feasibility gate returned
a qualified pass, and the qualification is disabling for the headline result: the reachable
portion of the dataset supports eight usable conditions, which cannot fund a task split with
more than one held-out test task, and a single test task cannot carry the paired interval
that every other result in this report carries. Restoring the remaining three quarters of the
images depends on a request form submitted to the dataset's authors, which is outside our
control and slow. Deliverable 9 is optional and conditional on deliverable 7, and is
deferred with it.

\paragraph{Deliverable 10, the few-shot face generation stress test.} Deferred by the plan's
own exclusion list, which directs that no further time be spent on it during this
submission cycle.

\paragraph{The dSprites intermediate stage.} Deferred by the same exclusion list. Its
intended role was to supply a task variable that cannot be read from a single target image,
and that role has been reassigned to the demographic experiment. We note that
Section~\ref{sec:a1} gives this decision more weight than it had when it was taken: the
advantage over target-only modelling disappears by twenty target images precisely because a
single target image already reveals the task, which is the condition dSprites was designed
to remove.

\paragraph{The severity sweep.} Deferred by the exclusion list unless required to interpret
a Stage A result. It was not required. The command-line flag remains wired and verified.

\paragraph{A sweep of the training-time source budget.} We considered whether the flat A4
curve might reflect what the transport map learned under a fixed training budget rather than
a property of the task. Testing that would require retraining at several source budgets. The
plan contains no such experiment: every mention of the source budget in it refers to the
evaluation script, and the plan's overall posture is that the implementation is frozen. It is
therefore recorded here as an open question and not run.

\paragraph{Unconstrained fine-tuning references.} Full fine-tuning of the network weights,
with and without abundant target data, is implemented and reachable but was not run this
stage. It is the measurement that would attribute the oracle gap reported in
Section~\ref{sec:a1} either to the task or to the basis and backbone, and we identify it as
the cheapest informative next step.

\section{Limitations}

Four limitations travel with every number in this report.

First, all CIFAR-100 results come from a \emph{single checkpoint} and a \emph{single seed}.
Every interval is therefore a within-run interval over episodes and contains no seed-to-seed
component. This matters most in Section~\ref{sec:a5}, where four separately trained models
are compared and run-to-run variation is the dominant nuisance term; the floor reported there
is estimated from four runs, not bounded.

Second, all readings are on the validation split. The test split is untouched and is
reserved for a single final number.

Third, absolute generative quality is unmeasured. No Fr\'echet or kernel Inception distance
is implemented in the project, and no real-against-real control establishes the floor of the
sliced Wasserstein distance at these sample sizes. The oracle row functions as a practical
floor, but the absolute level of any distance reported here should not be read as a distance
from zero. What this report establishes is that the coordinate controls generation, not that
the generated images are good.

Fourth, the two instruments used here disagree in at least one place where the disagreement
changes a conclusion, and the resolution in Section~\ref{sec:a1} rests on an argument about
what each instrument can see rather than on a third measurement.

\section{Conclusion}

The criterion fixed before the stage began is met, and met narrowly: knowing the source task
and the relation beats both relation-only transfer and target-only modelling when target
data are scarce, and the margin over target-only modelling is gone by twenty target images.
The positive control confirms that the mechanism is sound where the ground truth is known,
and locates the difficulty on images precisely: the well is shallow, not absent. Two
supporting results are negative in a useful way. Source evidence saturates below the smallest
budget we measured, so the method is cheap on the source side. And the non-linear transport
map is not distinguishable from a linear one, so the architecture can be simplified rather
than defended.

The single most informative next measurement is the unconstrained fine-tuning reference,
which would establish whether the remaining gap to the oracle belongs to the task or to the
architecture. The branch integration that unblocks deliverable 3 is the next most valuable,
and is a repository operation rather than an experiment.

\end{document}
